
%%=============================================================%%
%% METHODS  (CS theory relocated here)                         %%
%%=============================================================%%
\section{Methods}
\label{sec:method}

\subsection{Datasets}

We evaluate on two colorectal-cancer histopathology benchmarks.
\textit{NCT-CRC-HE-100K}~\cite{nct_dataset} contains 100{,}000 non-overlapping
$224\times224$ H\&E patches at 0.5\,$\mu$m/pixel across nine tissue classes
(ADI, BACK, DEB, LYM, MUC, MUS, NORM, STR, TUM). We use 5-fold cross-validation
with stratified splits (60\% train, 20\% validation, 20\% test per fold) and
report mean\,$\pm$\,s.d. across folds. \textit{CRC-VAL-HE-7K}~\cite{nct_dataset}
contains 7{,}180 patches from an independent cohort with different staining
protocols and scanners, used exclusively for cross-institution testing without
fine-tuning.

\subsection{Baselines and evaluation}

We compare DCA against 14 methods spanning standard CNNs (ResNet-18/50,
DenseNet-121, EfficientNet-B0), spatial-transformer methods (STN-Affine,
STN-TPS~\cite{jaderberg2015spatial}, Deformable
Convolution~\cite{dai2017deformable}), fixed geometric normalisation (CEM and
SEM~\cite{huang2024surface}), stain normalisation
(Macenko~\cite{macenko2009method}, StainGAN~\cite{shaban2019staingan}),
self-supervised pretraining (SimCLR~\cite{chen2020simple}), and pathology
foundation models (PLIP~\cite{huang2023plip}, UNI~\cite{chen2024uni},
CONCH~\cite{lu2024conch}) under both linear-probe and fine-tuned protocols. All
spatial-transformer and geometric methods use a ResNet-18 backbone for fair
comparison. We report top-1 accuracy, macro-F1, and Cohen's $\kappa$, and assess
significance with paired $t$-tests across folds with Bonferroni correction
($\dagger$: $p<0.05$; $\ddagger$: $p<0.01$).

\textbf{Implementation.} All models train for 50 epochs with AdamW
($\beta_1=0.9$, $\beta_2=0.999$, initial learning rate $10^{-3}$, weight decay
$10^{-4}$), cosine annealing with 5-epoch warmup, and batch size 32 on NVIDIA
A100 GPUs. Augmentation comprises random flips, rotation ($\pm15^\circ$), and
colour jitter ($\pm0.2$), disabled at validation and test. For DCA we set
$\lambda_{conf}=1.0$, $\lambda_{smooth}=0.1$, and $T=7$ (Table~\ref{tab:hyperparameter}).
Random seeds are fixed to \{42, 123, 456, 789, 1024\} for the five folds and
applied identically across methods. Foundation models use a $10\times$ lower
backbone learning rate for fine-tuning; SimCLR is pretrained for 200 epochs
(temperature 0.5) before supervised fine-tuning. Table~\ref{tab:comparison_prior}
summarises how DCA differs from prior geometric-transformation approaches.

\begin{table}[t]
\centering
\caption{Comparison of Geometric Transformation Approaches}
\label{tab:comparison_prior}
\begin{tabularx}{\textwidth}{Xcccc}
\hline
\textbf{Method} & \textbf{Learnable} & \textbf{Conformal} & \textbf{Diffeo.} & \textbf{End-to-End} \\
\hline
Standard STN~\cite{jaderberg2015spatial} & \checkmark & $\times$ & $\times$ & \checkmark \\
VoxelMorph~\cite{balakrishnan2019voxelmorph} & \checkmark & $\times$ & \checkmark & \checkmark \\
Fixed CEM~\cite{huang2024surface} & $\times$ & \checkmark & $\checkmark^*$ & $\times$ \\
\textbf{DCA (Ours)} & \checkmark & \checkmark & \checkmark & \checkmark \\
\hline
\end{tabularx}
\vspace{1mm}
\footnotesize{$^*$CEM produces conformal maps which are diffeomorphic under appropriate boundary conditions; however, the optimization is decoupled from classification.}
\end{table}


\subsection{Problem formulation}

Let $\Omega\subset\mathbb{R}^2$ denote the image domain and
$\mathcal{I}=L^2(\Omega,\mathbb{R}^3)$ the space of RGB images. Given a labelled
dataset $\mathcal{D}=\{(I_i,y_i)\}_{i=1}^N$ with $I_i\in\mathcal{I}$ and
$y_i\in\{1,\ldots,C\}$, we jointly learn a transformation network
$\mathcal{T}_\psi:\mathcal{I}\rightarrow\mathfrak{X}(\Omega)$ that maps each image
to a velocity field $v_i=\mathcal{T}_\psi(I_i)$ (in practice an $H\times W\times2$
tensor), from which a diffeomorphism $\phi_i=\exp(v_i)\in\text{Diff}^+(\Omega)$ is
obtained, and a classifier $f_\theta:\mathcal{I}\rightarrow\Delta^{C-1}$. The
transformed image is $\tilde{I}_i=I_i\circ\phi_i$ (backward warping via
differentiable bilinear interpolation). The learning problem is
\begin{equation}
\min_{\psi,\theta}\;\mathbb{E}_{(I,y)\sim\mathcal{D}}\left[\mathcal{L}_{cls}(f_\theta(\tilde{I}),y)+\lambda_{conf}\mathcal{E}_{conf}(\phi)+\lambda_{smooth}\mathcal{R}_{smooth}(v)\right],
\label{eq:objective}
\end{equation}
balancing classification accuracy, angle preservation, and deformation
smoothness.

\subsection{Diffeomorphic spatial transformer}

Rather than predicting $\phi$ directly, we predict a stationary velocity field
$v$ and obtain $\phi$ via the exponential map, $\phi=\exp(v)$
(Definition~\ref{def:exp_map}). The localisation network $g_\psi$ is a
U-Net~\cite{ronneberger2015unet} with four encoder blocks (channels
$[3,64,128,256,512]$) and a mirrored decoder with skip connections; the final
$1\times1$ convolution outputs the two velocity components and is initialised to
zero so that the transformation begins at identity. The exponential map is
approximated by scaling-and-squaring~\cite{arsigny2006logdiff}
(Algorithm~\ref{alg:scaling_squaring}), which guarantees a diffeomorphism for
bounded velocity magnitude (Proposition~\ref{prop:diffeomorphism}).

\begin{definition}[Exponential map]
\label{def:exp_map}
The exponential map $\exp:\mathfrak{X}(\Omega)\rightarrow\text{Diff}^+(\Omega)$ is
$\phi=\exp(v)=\lim_{T\rightarrow\infty}\left(\text{Id}+v/T\right)^T$, where
$\text{Id}$ is the identity transformation.
\end{definition}

\begin{algorithm}[t]
\caption{Scaling-and-Squaring for Diffeomorphic Integration}
\label{alg:scaling_squaring}
\begin{algorithmic}[1]
\REQUIRE Velocity field $v \in \mathbb{R}^{H \times W \times 2}$, number of steps $T \in \mathbb{N}$
\ENSURE Deformation field $\phi \approx \exp(v)$
\STATE $\phi^{(0)} \leftarrow v / 2^T$ \quad \COMMENT{$\phi^{(k)}$ represents displacement field}
\FOR{$k = 1$ to $T$}
    \STATE $\phi^{(k)}(x) \leftarrow \phi^{(k-1)}(x) + \phi^{(k-1)}(x + \phi^{(k-1)}(x))$ \quad \COMMENT{via bilinear interpolation}
\ENDFOR
\RETURN $\phi^{(T)}$
\end{algorithmic}
\end{algorithm}


\begin{proposition}[Diffeomorphism via scaling-and-squaring]
\label{prop:diffeomorphism}
Let $v:\Omega\rightarrow\mathbb{R}^2$ be a smooth stationary velocity field.
Algorithm~\ref{alg:scaling_squaring} produces $F(x)=x+\phi(x)$ approximating
$\exp(v)$ with error $O(2^{-T})$~\cite{arsigny2006logdiff}. If the scaled field
satisfies $\|v/2^T\|_{L^\infty}\leq\epsilon$ for sufficiently small $\epsilon$,
the result is a diffeomorphism, i.e.\ $\det(J_F(x))>0$ for all $x\in\Omega$.
\end{proposition}

In our experiments we verify empirically that $\min_x\det(J_F(x))>0.1$ across all
training and test samples and observe no folding artefacts.

\subsection{Conformal regularisation via differential geometry}

A smooth map is conformal if it preserves angles locally; in two dimensions this
is characterised by the Cauchy--Riemann equations (Definition~\ref{def:conformal}).
This property matters for histopathology because it maintains local structure
(nuclear spacing, glandular shape) while permitting normalisation of global
morphology.

\begin{definition}[Conformal map]
\label{def:conformal}
A smooth map $F:\Omega\rightarrow\Omega$ with Jacobian $J_F$ is conformal iff
$J_F$ is a scaled rotation at each point, equivalently iff it satisfies the
Cauchy--Riemann equations
$\partial F_1/\partial x=\partial F_2/\partial y$ and
$\partial F_1/\partial y=-\partial F_2/\partial x$.
\end{definition}

For $F(x)=x+\phi(x)$ we define the conformal energy as the $L^2$ deviation from
the Cauchy--Riemann equations,
\begin{equation}
\mathcal{E}_{conf}(\phi)=\frac{1}{|\Omega|}\int_\Omega\left[\left(\frac{\partial\phi_1}{\partial x}-\frac{\partial\phi_2}{\partial y}\right)^2+\left(\frac{\partial\phi_1}{\partial y}+\frac{\partial\phi_2}{\partial x}\right)^2\right]dx,
\label{eq:conformal_energy}
\end{equation}
discretised with central finite differences and replication padding. Strictly
conformal maps in 2D are highly restricted, so we appeal to \textit{quasiconformal}
mapping theory~\cite{ahlfors2006conformal}, which allows bounded angle distortion
measured by the Beltrami coefficient $\mu=\partial_{\bar z}F/\partial_z F$. A map
is conformal iff $\mu=0$, with maximal dilatation $K=(1+|\mu|)/(1-|\mu|)$
(Proposition~\ref{prop:beltrami}); minimising $\mathcal{E}_{conf}$ is equivalent
to minimising $\|\mu\|_{L^2}$, encouraging $K\approx1$. For typical learned
transformations we observe mean $|\mu|\approx0.15$ and $K\approx1.35$, i.e.\
quasiconformal maps with low distortion.

\begin{proposition}[Conformality and the Beltrami coefficient~\cite{ahlfors2006conformal}]
\label{prop:beltrami}
A smooth orientation-preserving map $F$ is conformal iff $\mu=0$ almost
everywhere. The maximal dilatation satisfies $K=(1+|\mu|)/(1-|\mu|)$, and
$\mathcal{E}_{conf}(\phi)=\frac{4}{|\Omega|}\int_\Omega|\mu|^2|\partial_zF|^2\,dx$
(see~\cite{lui2014brain}).
\end{proposition}

\subsection{End-to-end objective and training}

The full objective combines classification, conformal, and smoothness terms,
\begin{equation}
\mathcal{L}(\psi,\theta;I,y)=\mathcal{L}_{cls}(f_\theta(\tilde{I}),y)+\lambda_{conf}\mathcal{E}_{conf}(\phi)+\lambda_{smooth}\mathcal{R}_{smooth}(v),
\label{eq:full_loss}
\end{equation}
where $\mathcal{L}_{cls}$ is cross-entropy and
$\mathcal{R}_{smooth}(v)=\frac{1}{|\Omega|}\sum_{x}\|\nabla v(x)\|_F^2$ is a
total-variation regulariser on the velocity field, which favours piecewise-smooth
deformations while permitting sharp transitions at tissue boundaries. All
operations---scaling-and-squaring, bilinear warping, and conformal-energy
computation---are differentiable, enabling end-to-end optimisation by
backpropagation. The ResNet-18 classifier is initialised from ImageNet and the
U-Net final layer from zero, so training begins at the identity transformation.
Algorithm~\ref{alg:training} gives the complete procedure.

\begin{algorithm}[t]
\caption{DCA Training Procedure}
\label{alg:training}
\begin{algorithmic}[1]
\REQUIRE Dataset $\mathcal{D}$, batch size $B$, epochs $E$, initial learning rate $\alpha_0$
\ENSURE Trained parameters $\psi^*, \theta^*$
\STATE Initialize $g_\psi$ (U-Net): He init for conv layers, zero init for final $1 \times 1$ layer
\STATE Initialize $f_\theta$ (ResNet-18): ImageNet pretrained weights
\FOR{epoch $e = 1$ to $E$}
    \STATE Set learning rate $\alpha \leftarrow \alpha_0 \cdot 0.95^{\lfloor e/10 \rfloor}$
    \FOR{each mini-batch $\{(I_b, y_b)\}_{b=1}^B$}
        \STATE $v_b \leftarrow g_\psi(I_b)$ \quad \COMMENT{Predict velocity field}
        \STATE $\phi_b \leftarrow \text{ScalingSquaring}(v_b, T=7)$ \quad \COMMENT{Integrate to displacement}
        \STATE $\tilde{I}_b \leftarrow I_b \circ \phi_b$ \quad \COMMENT{Warp via bilinear interpolation}
        \STATE $\hat{y}_b \leftarrow f_\theta(\tilde{I}_b)$ \quad \COMMENT{Classify transformed image}
        \STATE Compute loss $\mathcal{L}(\psi, \theta; I_b, y_b)$ via Eq.~\eqref{eq:full_loss}
        \STATE Update $(\psi, \theta)$ via Adam with learning rate $\alpha$
        \STATE \textit{Optional:} Monitor $\min_x \det(J_{F_b}(x))$ to verify diffeomorphism
    \ENDFOR
\ENDFOR
\RETURN $(\psi^*, \theta^*)$
\end{algorithmic}
\end{algorithm}


\textbf{Computational cost.} The per-image forward pass is dominated by the two
network passes, giving $O(HW(D_U+D_R))$ complexity---identical in order to
standard CNN classification. Scaling-and-squaring adds only $\sim$2\,ms to the
$\sim$15\,ms total inference on an A100 GPU. In contrast to fixed conformal
methods~\cite{huang2024surface} that require $O(n^3)$ per-image optimisation
(30--60\,s per image), DCA amortises the geometric cost during training
(Table~\ref{tab:computational}), enabling real-time deployment.

\begin{table}[t]
\centering
\caption{Computational Cost Comparison (224$\times$224 input, A100 GPU)}
\label{tab:computational}
\begin{tabularx}{\textwidth}{Xccccc}
\hline
\textbf{Method} & \textbf{Params} & \textbf{FLOPs} & \textbf{Train} & \textbf{Infer.} & \textbf{Mem.} \\
 & (M) & (G) & (h) & (ms) & (GB) \\
\hline
ResNet-18 & 11.7 & 1.8 & 2.1 & 8.2 & 1.8 \\
ResNet-50 & 25.6 & 4.1 & 3.8 & 12.5 & 3.2 \\
STN-Affine & 12.1 & 2.0 & 2.4 & 9.5 & 2.0 \\
STN-TPS & 13.8 & 2.4 & 2.9 & 12.3 & 2.4 \\
Deform. Conv & 12.9 & 2.3 & 3.2 & 11.8 & 2.6 \\
CEM (online) & 11.7 & 1.8 & -- & 1,850 & 3.2 \\
CEM (precomp.) & 11.7 & 1.8 & -- & 8.2$^*$ & 1.8 \\
\textbf{DCA} & 19.2 & 3.1 & 3.5 & 11.1 & 2.6 \\
\hline
\end{tabularx}

\vspace{1mm}
\footnotesize{$^*$Excludes precomputation time (1,850ms/image offline).}
\end{table}
